\section{Introduction}
\label{sec:introduction}

The popularity of blockchains has undeniably changed the landscape of Byzantine fault tolerant (BFT) consensus.
Scalable protocols with multiple leaders, enabling transaction proposals to occur in parallel, are being pursued both in theory and in practice
~\cite{miller2016honey,crain2021red,yang2022dispersedledger, RCC, Mir, ISS, lev2019fairledger, 
 baird2016swirlds, danezis2018blockmania, gkagol2019aleph, danezis2022narwhal, spiegelman2022bullshark, keidar2022cordial, fino}. 
While these protocols demonstrate terrific throughput, they all have a similar ``recipe'': 
They allow parallel broadcasts to proceed for a while, and then wait for a leader to drive a commit decision on an entire ``wave'' of accumulated broadcasts. 
This introduces a rather large delay, as these protocols first wait for reliable broadcasts to be finalized and then start ordering them.

Traditional, non-BFT consensus protocols have benefited greatly from decoupling finalization and sequencing. 
Distributed shared logs~\cite{balakrishnan2012corfu, balakrishnan2013tango, wei2017vcorfu, balakrishnan2020virtual, 10.1145/3477132.3483544}
have shown that finalization and sequencing can be performed independently, and can greatly improve parallelism, throughput and latency. 
Instead of waiting for a consensus protocol to both finalize transactions and order them into a sequenced log, 
they assign slots in the log for transactions by a ``sequencer'' immediately upon request.
Sequenced transactions are then finalized into their assigned slots independently and in parallel, enabling the protocol to proceed as
quickly as the sequencer can assign slots.

In this work, \system, we apply the benefits of decoupling finalization from sequencing to the Byzantine setting. 
How do we build a sequencer in the Byzantine setting, though? Our first insight is remarkably simple:
consistent broadcast itself prevents sender equivocation, 
hence broadcasts can be finalized into pre-assigned slots so long as the uniqueness of slot numbers is enforced by the protocol.
Therefore, as in non-BFT settings, broadcasts carrying unique slot numbers can be finalized independently and in parallel, 
successfully growing a sequenced log.
Based on this insight, \system is built on a core broadcast primitive called \primitive, 
a novel variant of consistent broadcast guaranteeing unique slot numbering.
When there is no contention over slots numbers, \primitive broadcasts can be finalized in three network delays.

However, gaps in the log may occur in \system due to faulty nodes or network partitions. 
We use a fallback consensus decision to finalize gap slots. 
Harnessing recent advances in Byzantine fault tolerance on DAGs~\cite{spiegelman2022bullshark,fino},
we allow broadcasts to keep being finalized and drive consensus decisions with zero message overhead.
In this manner, while a gap in slot numbers may delay growing the committed log, 
all broadcasts have utility in contributing to throughput and system resources continue being fully utilized.

Combining these two components enables \system to decouple finalization and sequencing, providing a optimistically fast, high-throughput happy path with a high-thoughput fallback.
More specifically, \primitive allows blocks of transactions to be finalized into a sequenced log extremely fast in the common case. 
At the same time, using a DAG to drive consensus allows \system to maintain the high-throughput scalability of state-of-art DAG protocols (e.g.,~\cite{gkagol2019aleph, danezis2022narwhal, spiegelman2022bullshark}) under contention.
Note that this line of research work emphasizes separating transaction dissemination from block finalization to improve throughput. 
This is orthogonal yet compatible with the core design of \system.

As an added benefit, \system supports ``eager'' notifications which allow nodes to observe finalized slots in the log out of order.
Many use-cases can benefit from the eager finalization notifications, so that \system can support diverse applications on the same platform.

This paper presents a complete \system protocol, which is part of a larger ongoing system project. The goal of this paper is to
present the background on the key insights of \system, to formally describe the protocol and the key primitive \primitive,
and to argue correctness. We are in the process of building a complete system around \system~which
leverages its full performance and sequencing capabilities.